\chapter{Framework de Priorización Extendido}
\label{anexo:framework-priorizacion-extendido}

\setlength{\parindent}{0pt}

\section{Matriz de 18 escenarios de priorización}

La matriz presenta las combinaciones más comunes de valores altos (7--10), medios (4--6) y bajos (1--3) en cada dimensión, junto con la decisión típica. Sirve como guía inicial, no como regla definitiva; cada equipo desarrolla su propia interpretación basada en su experiencia acumulada.

\begingroup
\scriptsize
\renewcommand{\arraystretch}{1.4}
\setlength{\tabcolsep}{4pt}
\setlength{\LTleft}{0pt}
\setlength{\LTright}{0pt}

\begin{longtable}{
    >{\Centering\hspace{0pt}}p{0.8cm}     % Col 1: #
    >{\Centering\hspace{0pt}}p{1.5cm}     % Col 2: Valor
    >{\Centering\hspace{0pt}}p{1.5cm}     % Col 3: Técnica
    >{\Centering\hspace{0pt}}p{1.5cm}     % Col 4: Social
    >{\Centering\hspace{0pt}}p{2.2cm}     % Col 5: Prioridad
    >{\RaggedRight\hspace{0pt}}p{5.5cm}   % Col 6: Decisión típica
    }

    % --- 1. ENCABEZADO PARA LA PRIMERA PÁGINA ---
    \caption{Matriz de 18 escenarios de priorización} \label{tab:matriz-priorizacion}                                                                                   \\
    \toprule
    \textbf{\#} & \textbf{Valor} & \textbf{Técnica} & \textbf{Social} & \textbf{Prioridad}  & \textbf{Decisión típica}                                                  \\
    \midrule
    \endfirsthead

    % --- 2. ENCABEZADO PARA PÁGINAS SIGUIENTES ---
    \multicolumn{6}{c}{{\bfseries \tablename\ \thetable{} -- Continuación}}                                                                                             \\
    \toprule
    \textbf{\#} & \textbf{Valor} & \textbf{Técnica} & \textbf{Social} & \textbf{Prioridad}  & \textbf{Decisión típica}                                                  \\
    \midrule
    \endhead

    % --- 3. PIE PARA PÁGINAS INTERMEDIAS ---
    \midrule
    \multicolumn{6}{r}{\textit{Continúa en la siguiente página...}}                                                                                                     \\
    \endfoot

    % --- 4. PIE FINAL ---
    \bottomrule
    \endlastfoot

    1           & Alto           & Alto             & Bajo            & \textbf{Máxima}     & Entregar inmediatamente. El equipo sano puede absorber la complejidad.    \\
    2           & Alto           & Bajo             & Alto            & \textbf{Alta}       & Mitigar lo social primero. Un equipo en crisis no entrega calidad.        \\
    3           & Bajo           & Alto             & Alto            & \textbf{Alta}       & Resolver ambos en paralelo. La urgencia técnica y social no compiten.     \\
    4           & Alto           & Alto             & Alto            & \textbf{Escalar}    & Crisis total. Escalar al CTO para una decisión organizacional.            \\
    5           & Alto           & Medio            & Medio           & \textbf{Alta}       & Entregar con precaución. Monitorear la salud durante la implementación.   \\
    6           & Medio          & Alto             & Medio           & \textbf{Alta}       & Resolver lo técnico primero. La deuda técnica bloquea el progreso.        \\
    7           & Medio          & Medio            & Alto            & \textbf{Alta}       & Resolver lo social primero. Un equipo en crisis contamina todo lo demás.  \\
    8           & Alto           & Bajo             & Alto            & \textbf{Media-Alta} & Balancear: MVP de la \textit{feature} más intervención social.            \\
    9           & Bajo           & Alto             & Bajo            & \textbf{Media}      & Resolver lo técnico. Importante aunque no genere valor inmediato.         \\
    10          & Bajo           & Bajo             & Alto            & \textbf{Media}      & Resolver lo social. La salud del equipo es un prerrequisito.              \\
    11          & Alto           & Alto             & Bajo            & \textbf{Media-Alta} & Entregar. Equipo sano; aprovechar la capacidad disponible.                \\
    12          & Bajo           & Bajo             & Bajo            & \textbf{Baja}       & Posponer. No hay urgencias ni impacto inmediato.                          \\
    13          & Alto           & Bajo             & Bajo            & \textbf{Alta}       & Entregar. \textit{Feature} valiosa sin impedimentos relevantes.           \\
    14          & Bajo           & Alto             & Bajo            & \textbf{Media}      & Resolver lo técnico. Deuda urgente aunque el valor sea bajo.              \\
    15          & Bajo           & Bajo             & Alto            & \textbf{Media}      & Resolver lo social. El equipo necesita atención prioritaria.              \\
    16          & Medio          & Medio            & Medio           & \textbf{Media}      & Caso por caso. Todo es moderado; decidir según el contexto.               \\
    17          & Alto           & Bajo             & Alto            & \textbf{Alta}       & \textit{Feature} y social en paralelo. MVP mientras se atiende al equipo. \\
    18          & Alto           & Alto             & Bajo            & \textbf{Alta}       & Resolver lo técnico primero. Bug crítico bloqueante; equipo sano.         \\
\end{longtable}
\endgroup

\section{Protocolo de decisión colaborativa}

La priorización ocurre durante la \textit{Social Sprint Planning} (15--20~min, antes de la estimación detallada). Ningún rol tiene autoridad absoluta; la decisión emerge del diálogo.

\subsection{Participantes y autoridad}

\begingroup
\scriptsize
\renewcommand{\arraystretch}{1.4}
\setlength{\tabcolsep}{4pt}
\setlength{\LTleft}{0pt}
\setlength{\LTright}{0pt}

\begin{longtable}{
    >{\RaggedRight\hspace{0pt}}p{3.0cm}   % Col 1: Rol
    >{\RaggedRight\hspace{0pt}}p{3.0cm}   % Col 2: Dimensión
    >{\RaggedRight\hspace{0pt}}p{8.0cm}   % Col 3: Autoridad
    }

    % --- 1. ENCABEZADO PARA LA PRIMERA PÁGINA ---
    \caption{Autoridad de cada rol en la priorización} \label{tab:autoridad-roles}                         \\
    \toprule
    \textbf{Rol}  & \textbf{Dimensión} & \textbf{Autoridad}                                                \\
    \midrule
    \endfirsthead

    % --- 2. ENCABEZADO PARA PÁGINAS SIGUIENTES ---
    \multicolumn{3}{c}{{\bfseries \tablename\ \thetable{} -- Continuación}}                                \\
    \toprule
    \textbf{Rol}  & \textbf{Dimensión} & \textbf{Autoridad}                                                \\
    \midrule
    \endhead

    % --- 3. PIE PARA PÁGINAS INTERMEDIAS ---
    \midrule
    \multicolumn{3}{r}{\textit{Continúa en la siguiente página...}}                                        \\
    \endfoot

    % --- 4. PIE FINAL ---
    \bottomrule
    \endlastfoot

    Product Owner & Valor de Negocio   & Decide \textbf{qué} entregar al cliente.                          \\
    Scrum Master  & Proceso            & Decide \textbf{cómo} trabajar. Protege el marco Scrum.            \\
    Social Guide  & Salud Social       & Decide \textbf{cuándo} el equipo puede comprometerse.             \\
    Dev Team      & Urgencia Técnica   & Decide \textbf{cuánto} es factible. Advierte sobre deuda técnica. \\
\end{longtable}
\endgroup

\subsection{Protocolo de priorización paso a paso}

\begin{description}[nosep, leftmargin=1.2cm, style=nextline]
    \item[Paso 1: Estado del equipo (5 min)] El Social Guide comparte el \textit{Health Score}, las dimensiones problemáticas, los \gls{community-smells} activos, la tendencia y la capacidad estimada (Tabla~\ref{tab:diagnostico-social-guide}). Si el HS es crítico, los pasos siguientes se ajustan.

    \item[Paso 2: \textit{Features} priorizadas (5 min)] El PO presenta los ítems de mayor valor (Tabla~\ref{tab:features-priorizadas-valor}), proponiendo más de los que caben en el \textit{Sprint} para permitir flexibilidad.

    \item[Paso 3: Ítems sociales (5 min)] Si el HS lo indica, el SG propone intervenciones con estimación en \textit{Social Story Points} (Tabla~\ref{tab:propuesta-social-guide}). Si el HS es saludable ($\geq 8$), no se requieren ítems sociales.

    \item[Paso 4: Discusión colaborativa (5--10 min)] Los cuatro roles dialogan para llegar a un \textit{Sprint Backlog} consensuado (Tabla~\ref{tab:decision-sprint-dialogo}).
\end{description}

\begingroup
\scriptsize
\renewcommand{\arraystretch}{1.4}
\setlength{\tabcolsep}{4pt}
\setlength{\LTleft}{0pt}
\setlength{\LTright}{0pt}

\begin{longtable}{
    >{\RaggedRight\hspace{0pt}}p{4.5cm}   % Col 1: Indicador
    >{\RaggedRight\hspace{0pt}}p{9.0cm}   % Col 2: Estado
    }

    % --- 1. ENCABEZADO PARA LA PRIMERA PÁGINA ---
    \caption{Ejemplo de diagnóstico del Social Guide} \label{tab:diagnostico-social-guide}  \\
    \toprule
    \textbf{Indicador}             & \textbf{Estado}                                        \\
    \midrule
    \endfirsthead

    % --- 2. ENCABEZADO PARA PÁGINAS SIGUIENTES ---
    \multicolumn{2}{c}{{\bfseries \tablename\ \thetable{} -- Continuación}}                 \\
    \toprule
    \textbf{Indicador}             & \textbf{Estado}                                        \\
    \midrule
    \endhead

    % --- 3. PIE PARA PÁGINAS INTERMEDIAS ---
    \midrule
    \multicolumn{2}{r}{\textit{Continúa en la siguiente página...}}                         \\
    \endfoot

    % --- 4. PIE FINAL ---
    \bottomrule
    \endlastfoot

    \textit{Health Score} actual   & 6.2/10 (zona: en riesgo)                               \\
    Dimensiones bajas              & Coraje (4.5), Foco (5.1)                               \\
    \gls{community-smells} activos & \textit{Radio Silence} (moderado)                      \\
    Tendencia                      & Mejorando (5.1 en el \textit{Sprint} anterior)         \\
    Capacidad estimada             & $\approx$80~\% de la normal (32 SP vs.\ 40 SP típicos) \\
\end{longtable}
\endgroup

\begingroup
\scriptsize
\renewcommand{\arraystretch}{1.4}
\setlength{\tabcolsep}{4pt}
\setlength{\LTleft}{0pt}
\setlength{\LTright}{0pt}

\begin{longtable}{
    >{\Centering\hspace{0pt}}p{1.0cm}     % Col 1: #
    >{\RaggedRight\hspace{0pt}}p{9.0cm}   % Col 2: Feature
    >{\Centering\hspace{0pt}}p{2.5cm}     % Col 3: Valor
    }

    % --- 1. ENCABEZADO PARA LA PRIMERA PÁGINA ---
    \caption{Ejemplo de \textit{features} priorizadas por valor de negocio (PO)} \label{tab:features-priorizadas-valor} \\
    \toprule
    \textbf{\#} & \textbf{\textit{Feature}}       & \textbf{Valor}                                                      \\
    \midrule
    \endfirsthead

    % --- 2. ENCABEZADO PARA PÁGINAS SIGUIENTES ---
    \multicolumn{3}{c}{{\bfseries \tablename\ \thetable{} -- Continuación}}                                             \\
    \toprule
    \textbf{\#} & \textbf{\textit{Feature}}       & \textbf{Valor}                                                      \\
    \midrule
    \endhead

    % --- 3. PIE PARA PÁGINAS INTERMEDIAS ---
    \midrule
    \multicolumn{3}{r}{\textit{Continúa en la siguiente página...}}                                                     \\
    \endfoot

    % --- 4. PIE FINAL ---
    \bottomrule
    \endlastfoot

    1           & Integración con API de Pagos    & 9/10                                                                \\
    2           & \textit{Dashboard} de Analytics & 7/10                                                                \\
    3           & Notificaciones \textit{push}    & 8/10                                                                \\
    4           & Mejora de UX en \textit{login}  & 5/10                                                                \\
    5           & Exportar reportes a PDF         & 6/10                                                                \\
\end{longtable}
\endgroup

\begingroup
\scriptsize
\renewcommand{\arraystretch}{1.4}
\setlength{\tabcolsep}{4pt}
\setlength{\LTleft}{0pt}
\setlength{\LTright}{0pt}

\begin{longtable}{
    >{\Centering\hspace{0pt}}p{1.0cm}     % Col 1: #
    >{\RaggedRight\hspace{0pt}}p{8.5cm}   % Col 2: Intervención propuesta
    >{\Centering\hspace{0pt}}p{3.0cm}     % Col 3: Esfuerzo (SSP)
    }

    % --- 1. ENCABEZADO PARA LA PRIMERA PÁGINA ---
    \caption{Ejemplo de propuesta de intervención del Social Guide (HS 6.2)} \label{tab:propuesta-social-guide} \\
    \toprule
    \textbf{\#} & \textbf{Intervención propuesta}                        & \textbf{Esfuerzo (SSP)}              \\
    \midrule
    \endfirsthead

    % --- 2. ENCABEZADO PARA PÁGINAS SIGUIENTES ---
    \multicolumn{3}{c}{{\bfseries \tablename\ \thetable{} -- Continuación}}                                     \\
    \toprule
    \textbf{\#} & \textbf{Intervención propuesta}                        & \textbf{Esfuerzo (SSP)}              \\
    \midrule
    \endhead

    % --- 3. PIE PARA PÁGINAS INTERMEDIAS ---
    \midrule
    \multicolumn{3}{r}{\textit{Continúa en la siguiente página...}}                                             \\
    \endfoot

    % --- 4. PIE FINAL ---
    \bottomrule
    \endlastfoot

    1           & Mediación Dev--QA para resolver \textit{Radio Silence} & 5                                    \\
    2           & Establecer protocolo de \textit{daily} asíncrona       & 3                                    \\
\end{longtable}
\endgroup

\noindent\textbf{Ejemplo de resultado del diálogo:} con capacidad de 32~SP, el equipo acuerda incluir API de Pagos (13~SP), Mediación Dev--QA (5~SSP), \textit{Dashboard} Analytics MVP (8~SP) y \textit{daily} asíncrona (3~SSP) = 29 puntos. Las notificaciones \textit{push} quedan para el siguiente \textit{Sprint}. El SG proyecta una mejora del HS hasta 7.0--7.5 si se abordan las intervenciones sociales.

\begingroup
\scriptsize
\renewcommand{\arraystretch}{1.4}
\setlength{\tabcolsep}{4pt}
\setlength{\LTleft}{0pt}
\setlength{\LTright}{0pt}

\begin{longtable}{
    >{\RaggedRight\hspace{0pt}}p{7.0cm}   % Col 1: Ítem comprometido
    >{\Centering\hspace{0pt}}p{3.0cm}     % Col 2: Tipo
    >{\Centering\hspace{0pt}}p{3.0cm}     % Col 3: Esfuerzo
    }

    % --- 1. ENCABEZADO PARA LA PRIMERA PÁGINA ---
    \caption{Compromiso final del Sprint tras diálogo de priorización} \label{tab:decision-sprint-dialogo} \\
    \toprule
    \textbf{Ítem comprometido}            & \textbf{Tipo}      & \textbf{Esfuerzo}                         \\
    \midrule
    \endfirsthead

    % --- 2. ENCABEZADO PARA PÁGINAS SIGUIENTES ---
    \multicolumn{3}{c}{{\bfseries \tablename\ \thetable{} -- Continuación}}                                \\
    \toprule
    \textbf{Ítem comprometido}            & \textbf{Tipo}      & \textbf{Esfuerzo}                         \\
    \midrule
    \endhead

    % --- 3. PIE PARA PÁGINAS INTERMEDIAS ---
    \midrule
    \multicolumn{3}{r}{\textit{Continúa en la siguiente página...}}                                        \\
    \endfoot

    % --- 4. PIE FINAL ---
    \bottomrule
    \endlastfoot

    Integración con API de Pagos          & Técnico            & 13 SP                                     \\
    Mediación Dev--QA                     & Social             & 5 SSP                                     \\
    \textit{Dashboard} de Analytics (MVP) & Técnico            & 8 SP                                      \\
    \textit{Daily} asíncrona              & Social             & 3 SSP                                     \\
    \midrule
    \multicolumn{2}{r}{\textbf{Total}}    & \textbf{29 puntos}                                             \\
\end{longtable}
\endgroup

\subsection{Conversión entre SSP y SP}

SocialScrum establece una equivalencia aproximada: $1\;\text{SSP} \approx 1\;\text{SP}$. Ambos tipos de ítems consumen tiempo y atención del equipo y pueden sumarse para calcular la capacidad total del \textit{Sprint}. Cada equipo calibra esta equivalencia según su contexto. La tabla~\ref{tab:decision-sprint-dialogo} ilustra un ejemplo de compromiso final que combina ambos tipos de puntos.

\subsection{Protocolo de discrepancia}

Cuando PO, SG y SM no logran consenso, se activa un protocolo de resolución estructurado:

\begin{enumerate}[nosep]
    \item \textbf{Argumentación (9 min):} cada rol dispone de 3~min para presentar su posición con evidencia (valor de negocio, estado del \textit{Health Score}, perspectiva de proceso).
    \item \textbf{Input del Dev Team (5 min):} el equipo expresa su perspectiva sobre factibilidad y riesgos.
    \item \textbf{Votación informal:} tres opciones --- (A)~priorizar \textit{feature} (aceptar riesgo social), (B)~priorizar salud social (posponer \textit{feature}), (C)~híbrido (MVP + mitigación social).
    \item \textbf{Decisión o escalamiento:} si hay mayoría clara, se implementa. Si no, se escala al CTO/VP Engineering con un resumen del dilema. La \textit{Planning} continúa con ítems de consenso mientras se espera la decisión (máx.\ 24~h).
\end{enumerate}

\section{Ejemplo completo: Equipo Alpha (Sprint 14)}

\begingroup
\scriptsize
\renewcommand{\arraystretch}{1.4}
\setlength{\tabcolsep}{4pt}
\setlength{\LTleft}{0pt}
\setlength{\LTright}{0pt}

\begin{longtable}{
    >{\RaggedRight\hspace{0pt}}p{4.5cm}   % Col 1: Indicador
    >{\RaggedRight\hspace{0pt}}p{9.0cm}   % Col 2: Estado
    }

    % --- 1. ENCABEZADO PARA LA PRIMERA PÁGINA ---
    \caption{Diagnóstico social inicial del equipo Alpha} \label{tab:diagnostico-alpha}         \\
    \toprule
    \textbf{Indicador}        & \textbf{Estado}                                                 \\
    \midrule
    \endfirsthead

    % --- 2. ENCABEZADO PARA PÁGINAS SIGUIENTES ---
    \multicolumn{2}{c}{{\bfseries \tablename\ \thetable{} -- Continuación}}                     \\
    \toprule
    \textbf{Indicador}        & \textbf{Estado}                                                 \\
    \midrule
    \endhead

    % --- 3. PIE PARA PÁGINAS INTERMEDIAS ---
    \midrule
    \multicolumn{2}{r}{\textit{Continúa en la siguiente página...}}                             \\
    \endfoot

    % --- 4. PIE FINAL ---
    \bottomrule
    \endlastfoot

    \textit{Health Score}     & 6.8/10 (zona en riesgo, límite superior)                        \\
    Dimensiones               & Apertura 7.2, Compromiso 7.0, Respeto 7.5, Foco 5.8, Coraje 6.5 \\
    \textit{Community smells} & \textit{Context Switching} (moderado)                           \\
    Tendencia                 & Estable (6.7 en el \textit{Sprint} anterior)                    \\
    Capacidad estimada        & 90~\% $\rightarrow$ 36 SP de 40 típicos                         \\
\end{longtable}
\endgroup

Tras el diálogo de priorización, el PO presentó 5 \textit{features} técnicas (39~SP totales) y el SG propuso una intervención social (\textit{Focus Fridays}, 3~SSP) para abordar el \textit{Context Switching} y mejorar la dimensión Foco. Con una capacidad de 36~SP, el equipo acordó el siguiente \textit{Sprint Backlog}, posponiendo el refactor de reportes (8~SP) al Sprint~15.

\begingroup
\scriptsize
\renewcommand{\arraystretch}{1.4}
\setlength{\tabcolsep}{4pt}
\setlength{\LTleft}{0pt}
\setlength{\LTright}{0pt}

\begin{longtable}{
    >{\RaggedRight\hspace{0pt}}p{7.0cm}   % Col 1: Ítem
    >{\Centering\hspace{0pt}}p{3.0cm}     % Col 2: Tipo
    >{\Centering\hspace{0pt}}p{3.0cm}     % Col 3: Puntos
    }

    % --- 1. ENCABEZADO PARA LA PRIMERA PÁGINA ---
    \caption{Sprint Backlog final --- Equipo Alpha} \label{tab:sprint-backlog-alpha}    \\
    \toprule
    \textbf{Ítem}                                   & \textbf{Tipo}   & \textbf{Puntos} \\
    \midrule
    \endfirsthead

    % --- 2. ENCABEZADO PARA PÁGINAS SIGUIENTES ---
    \multicolumn{3}{c}{{\bfseries \tablename\ \thetable{} -- Continuación}}             \\
    \toprule
    \textbf{Ítem}                                   & \textbf{Tipo}   & \textbf{Puntos} \\
    \midrule
    \endhead

    % --- 3. PIE PARA PÁGINAS INTERMEDIAS ---
    \midrule
    \multicolumn{3}{r}{\textit{Continúa en la siguiente página...}}                     \\
    \endfoot

    % --- 4. PIE FINAL ---
    \bottomrule
    \endlastfoot

    Migración Auth                                  & Técnico         & 13 SP           \\
    API Partners                                    & Técnico         & 8 SP            \\
    Optimización de \textit{queries}                & Técnico         & 5 SP            \\
    Mejora \textit{dashboard} admin                 & Técnico         & 5 SP            \\
    \textit{Focus Fridays}                          & Social          & 3 SSP           \\
    \midrule
    \multicolumn{2}{r}{\textbf{Total comprometido}} & \textbf{34 pts}                   \\
    \multicolumn{2}{r}{\textbf{Capacidad}}          & \textbf{36 SP}                    \\
    \multicolumn{2}{r}{\textbf{Margen}}             & \textbf{2 SP}                     \\
\end{longtable}
\endgroup

\noindent\textbf{Sprint Goal}: Habilitar microservicios y mejorar el foco del equipo.